% ==========================================================
% STANDARDIZED PREAMBLE (STATIC CORE) — DO NOT EDIT THIS PART
% ==========================================================

\documentclass[11pt,a4paper]{article}

\usepackage{iftex}
\ifPDFTeX
  \usepackage[utf8]{inputenc}
  \usepackage[T1]{fontenc}
  \usepackage{stix2}
\else
  \usepackage{fontspec}
  \usepackage{unicode-math}
  \setmainfont{STIX Two Text}
  \setmathfont{STIX Two Math}
\fi

\usepackage[margin=2.5cm]{geometry}
\usepackage{parskip}
\usepackage{microtype}
\usepackage{enumitem}

\usepackage{amsmath, amssymb}

\usepackage{tocloft}
\setlength{\cftbeforesecskip}{2pt}
\setlength{\cftbeforesubsecskip}{1pt}

\widowpenalty=10000
\clubpenalty=10000

\usepackage{xcolor}
\definecolor{myblue}{RGB}{0, 51, 153}

\usepackage{sectsty}
\partfont{\color{myblue}}
\chapterfont{\color{myblue}}
\sectionfont{\color{myblue}}
\subsectionfont{\color{myblue}}
\subsubsectionfont{\color{myblue}}

\usepackage{hyperref}

\renewcommand{\cftsecfont}{\bfseries\color{myblue}}
\renewcommand{\cftsubsecfont}{\color{myblue}}
\renewcommand{\cftsubsubsecfont}{\color{myblue}}
\renewcommand{\cftsecpagefont}{\bfseries}
\renewcommand{\cftsubsecpagefont}{}
\renewcommand{\cftsubsubsecpagefont}{}

\hypersetup{
  colorlinks=true,
  linkcolor=myblue,
  citecolor=myblue,
  urlcolor=myblue
}

\usepackage{titlesec}
\titlespacing{\section}{0pt}{12pt}{6pt}
\titlespacing{\subsection}{0pt}{10pt}{5pt}
\titlespacing*{\section}{0pt}{2.5ex plus 1ex minus .2ex}{1.5ex plus .2ex}
\titlespacing*{\subsection}{0pt}{2ex plus 0.8ex minus .2ex}{1.2ex plus .2ex}
\titlespacing*{\subsubsection}{0pt}{1.8ex plus 0.6ex minus .2ex}{1ex plus .2ex}

\makeatletter
\newcommand{\StartBlueDoc}[3][]{%
  \title{\vspace{-2em}\textbf{#2}%
    \if\relax\detokenize{#3}\relax\else\\[2pt]{\large #3}\fi
    \vspace{+0.6em}}
  \author{#1}
  \date{\vspace{-0.6em}\today}
  \begin{document}
  \maketitle
  \tableofcontents
  \vspace{1em}\hrule\vspace{1em}
}
\makeatother

% ==========================================================
% CUSTOMIZATION ZONE
% ==========================================================


\newcommand{\DocAuthor}{}
\newcommand{\DocTitle}{Documenting Research Software}
\newcommand{\DocSubtitle}{}

\setcounter{tocdepth}{2}

\usepackage{booktabs}
\usepackage{tabularx}
\newcolumntype{L}{>{\raggedright\arraybackslash}X}
\newcolumntype{T}{>{\raggedright\arraybackslash}p{0.36\textwidth}}

% Header rows of the table "At a glance": a section, with its number and its
% title.
\newcommand{\GlanceSection}[1]{%
  \multicolumn{3}{@{}l}{\textbf{\ref{#1}\quad\nameref{#1}}}}

% Cite a principle of the other documents by its title. A \ref cannot
% reach into another PDF, so the title is the reference.
\newcommand{\DesignRef}[1]{\emph{Design Principles}, ``#1''}
\newcommand{\TestingRef}[1]{\emph{Testing Research Software}, ``#1''}
\newcommand{\TypesRef}[1]{\emph{Types of Research Software Tests},
  ``#1''}

% ==================
% Start the document
% ==================
\StartBlueDoc[\DocAuthor]{\DocTitle}{\DocSubtitle}



% ==========================================================
\section*{About this guide}
\phantomsection
\addcontentsline{toc}{section}{About this guide}
\label{sec:about}
% ==========================================================

This guide is about documenting research software: organizing, writing, and
maintaining what users and maintainers read besides the code. Together with
\emph{Design Principles for Research Software}, cited below as \emph{Design
Principles}, and \emph{Testing Research Software}, it forms one set of
guidelines for all code written for a research purpose, in academia or in
industry, in any programming language and scientific field. It is meant for
looking things up while working, not for reading from cover to cover.

Users act on documentation. They choose parameters by it, interpret their
results by it, and copy its description of the method into the methods
sections of their publications. Missing documentation is easy to notice;
wrong documentation is not. A description of last year's default, or of the
method as a paper published it before the code changed, gives no error, and
its consequences show up far away, in a result reported with the wrong
method. Two ideas run through this guide. Readers come to documentation with
different needs, and each kind of document serves one
(Section~\ref{sec:what}). And documentation stays true only if it lives,
changes, and is checked together with the code (Section~\ref{sec:true}).
Section~\ref{sec:writing} is about writing the reference, the comments, the
changelog, and the citation.

What has to be documented belongs to \emph{Design Principles}, because it is
part of what the software promises. Its principles require, among other
things, the method with its assumptions, defaults, and references
(\DesignRef{Document the method, not the implementation}), the policy for
missing values (\DesignRef{Mind the gap, don't fill it}), the time and
memory complexity (\DesignRef{Size matters}), a changelog (\DesignRef{Design
stable APIs}), and citation metadata (\DesignRef{Make every result traceable
to a version}). This guide says how such documentation is organized,
written, and kept true. Documentation explains correct use, but it does not
enforce it: users who never read it are protected only by the checks of the
code (\DesignRef{Garbage in, useful error out}).



\subsection*{Normative language and the ladder}
\phantomsection
\addcontentsline{toc}{subsection}{Normative language and the ladder}
\label{sec:normative}

The words \emph{must}, \emph{should}, and \emph{may}, and the plain
imperatives ``prefer'' and ``avoid'', carry the meaning defined in the
subsection ``Normative language'' of \emph{Design Principles}. The rules of
this guide apply along the same ``Ladder of Research Software'': the
\emph{must}s are part of the floor from the second rung on, and what rises
with the rung is how much documentation there is and for whom it is written
(Subsection~\ref{sec:grow}). When rules conflict, the priorities in
\DesignRef{Get your priorities right} decide. Each rule is stated in one of
the documents only.



\subsection*{Terminology}
\phantomsection
\addcontentsline{toc}{subsection}{Terminology}
\label{sec:terminology}

The terms defined in \emph{Design Principles}, such as package, library code,
workflow code, entry point, user-facing function, host stack, and caller,
carry the same meaning here, and so do the terms of testing defined in
\emph{Types of Research Software Tests}, such as test, oracle, and snapshot.
This guide adds the four kinds of documentation that
Subsection~\ref{sec:diataxis} keeps apart:

\begin{description}
  \item[Tutorial] A lesson that takes a newcomer through a complete task, step
    by step, so that they learn to use the software.
  \item[How-to guide] Directions for one task that a user who already knows
    the software wants to get done.
  \item[Reference] The description of every entry point, option, and file
    format, made for looking things up.
  \item[Explanation] The discussion of a topic, such as the method, its
    assumptions, and the reasons for a design, made for understanding.
\end{description}



% The checklist gets a page of its own, so that it can be printed.
\clearpage
% ==========================================================
\section*{At a glance}
\phantomsection
\addcontentsline{toc}{section}{At a glance}
% ==========================================================

Every principle of this guide in one line. The number and the title link to
the principle.

{\small
\begin{tabularx}{\textwidth}{@{}lTL@{}}
  \toprule
  \GlanceSection{sec:what} \\
  \midrule
  \ref{sec:diataxis} & \nameref{sec:diataxis}
    & Give each document one purpose, and provide the kinds that the readers
      need. \\
  \ref{sec:grow} & \nameref{sec:grow}
    & Write what the rung demands, from a README on the second rung to
      documentation that users can work from alone. \\
  \ref{sec:readme} & \nameref{sec:readme}
    & Say what the software does, its status, how to install, run, and cite
      it, and where to go next. \\
  \ref{sec:paper} & \nameref{sec:paper}
    & Cite the paper, but document the method as the current version
      implements it, and every difference. \\
  \addlinespace
  \GlanceSection{sec:writing} \\
  \midrule
  \ref{sec:reference} & \nameref{sec:reference}
    & Write the reference in the source, next to the code it describes, and
      generate it from there. \\
  \ref{sec:workflow-docs} & \nameref{sec:workflow-docs}
    & Document every configuration key, every input and output file, and
      the resources a run needs. \\
  \ref{sec:comments} & \nameref{sec:comments}
    & Explain reasons, sources, and conventions in comments, and record
      design decisions in the repository. \\
  \ref{sec:changelog} & \nameref{sec:changelog}
    & Write the changelog for users, by release and kind of change, with
      changes of results set apart. \\
  \ref{sec:cite} & \nameref{sec:cite}
    & Say how to cite the software, the method, and the tools it runs, in
      every place users look. \\
  \addlinespace
  \GlanceSection{sec:true} \\
  \midrule
  \ref{sec:docs-as-code} & \nameref{sec:docs-as-code}
    & Keep the documentation in the repository, change it with the code,
      build it on every change, and publish it per version. \\
  \ref{sec:examples} & \nameref{sec:examples}
    & Run every example on every change, and check the output it shows. \\
  \bottomrule
\end{tabularx}}



% ==========================================================
\section{What to document}
\label{sec:what}
% ==========================================================

\subsection{Di\'ataxis: tutorials, how-to guides, reference, explanation}
\label{sec:diataxis}

Readers come to documentation with one of four needs: to learn the
software, to get a task done, to look something up, or to understand. The
Di\'ataxis framework assigns each need a kind of document: a tutorial, a
how-to guide, a reference, and an explanation (\hyperref[sec:terminology]{Terminology}).
Each document should serve one of these needs. A document that serves
several serves none well: a tutorial that stops to list every option loses
the newcomer, and a reference interrupted by derivations makes a reader who
came to look something up search for it.

In research software, the four kinds take typical forms:

\begin{itemize}
  \item \textbf{Tutorial.} A complete analysis on a small dataset that ships
    with the package or is generated, from the input to the result, with the
    output shown at each step.
  \item \textbf{How-to guide.} Directions for a task such as running on a
    cluster, using one's own reference genome or grid, reading another file
    format, or reproducing the analysis of a paper.
  \item \textbf{Reference.} The description of every user-facing function,
    command, configuration key, and file format
    (Subsections~\ref{sec:reference} and~\ref{sec:workflow-docs}).
  \item \textbf{Explanation.} The method: its derivation and assumptions,
    when it applies and when it does not, how it compares with the
    alternatives, and how it was validated
    (Subsection~\ref{sec:paper}).
\end{itemize}

Keep them apart, and link them. A tutorial links to the reference for the
functions it uses, and the reference of a function links to the explanation
of its method instead of repeating it (\DesignRef{DRY: Don't repeat
yourself}). Which of the four a package needs, and how much of each, depends
on its rung (Subsection~\ref{sec:grow}).



\subsection{Grow the documentation with the software}
\label{sec:grow}

Documentation costs effort to write and more to keep true, and it should
grow with the rung of the software on the ``Ladder of Research Software''
(\DesignRef{Pay as you grow}). Write the documentation that the readers of
the current rung need, and add the rest when the software climbs, before new
readers arrive.

On the first rung, nothing is required, but a few lines on how to run the
code and what it assumes help its author a year later. From the second rung
on, the floor of \emph{Design Principles} already requires much of the
reference: the method, the policy for missing values, and the other parts
listed in \hyperref[sec:about]{About this guide}; a README
(Subsection~\ref{sec:readme}) makes it usable by colleagues who can still
ask the authors. From the third rung on, users cannot ask, and the
ladder asks for documentation that they can work from alone: a complete
reference, at least one tutorial, how-to guides for the common tasks, and documentation published
for each version (Subsection~\ref{sec:docs-as-code}). On the fourth rung,
the explanation of the method, its limits, and its validation should be as
thorough as the cost of a wrong result demands.

Documentation that has gone out of date is worse than none, because readers
trust it. Do not write more than can be kept true: a short README that is
right serves users better than a long manual that describes an earlier
version.



\subsection{The README is the front door}
\label{sec:readme}

Every package from the second rung on should have a README, and it should
answer, in this order, the questions of someone who has just found the
package:

\begin{itemize}
  \item What it does, in one or two sentences, and for which field and data.
  \item Its status: the rung, for example whether it is maintained for others
    or shared as is, and the stability commitment (\DesignRef{Design stable
    APIs}). Where the package deviates from a \emph{should} because of its
    rung, the subsection ``Normative language'' of \emph{Design Principles}
    asks for the reason to be stated once for the whole package, and the
    README is the place.
  \item How to install it, including the supported versions of the language
    and of the host stack.
  \item A minimal example: the shortest complete use, on data that the
    reader has or the package provides.
  \item How to cite it (Subsection~\ref{sec:cite}), its license, and how to
    get help and report bugs.
  \item Where the full documentation is.
\end{itemize}

Keep the README short. It is the front door, not the house: it points to the
tutorials and the reference rather than repeating them, because every
repetition is a copy that can drift (\DesignRef{DRY: Don't repeat
yourself}). The minimal example is an example like any other and runs on
every change (Subsection~\ref{sec:examples}).



\subsection{The paper is not the manual}
\label{sec:paper}

Much research software comes with a paper that describes its method, and
the paper should be cited from the documentation (\DesignRef{Document the
method, not the implementation}). It cannot replace the documentation. A
paper describes the method as it was at the time of publication, for
reviewers rather than users, and it often leaves out what users need: the
defaults, the treatment of edge cases, and the options added later. The
paper is fixed, and the software moves on.

The documentation should therefore describe the method as the current
version implements it. \emph{Design Principles} requires it to state every
difference from the literature it cites (\DesignRef{Document the method, not
the implementation}), such as a changed default, a corrected formula, an
approximation, or an extension. Give each difference with the version in
which it appeared: users who cite the paper for their results need it to
report what was actually computed (\DesignRef{Make every result traceable to
a version}). Where the
method of a paper is meant to be reproduced exactly, say which version of
the software reproduces it, and keep that version available.



% ==========================================================
\section{Writing documentation}
\label{sec:writing}
% ==========================================================

\subsection{The reference lives next to the code}
\label{sec:reference}

Write the reference of each user-facing function in the source, next to its
code, in the host stack's format for documentation comments, such as
docstrings in Python, roxygen comments in R, or docstrings in Julia, and
generate the published reference from there with the host stack's tools,
such as Sphinx, pkgdown, or Documenter. A reference written in a separate
place drifts from the code it describes, and nobody notices; one written
next to the code is in view whenever the code changes, and in the same
change (Subsection~\ref{sec:docs-as-code}).

The reference of a user-facing function should state what it computes, each
parameter with its type, shape, unit, and default, the return value with
the same, the errors and warnings it raises, and a short example
(Subsection~\ref{sec:examples}). For library code, the meaning of each
parameter matters more than its type, which the signature often shows
already: not ``tolerance, a float'', but which error it bounds, whether it
is relative or absolute, and what happens when it is not met.

Where the tools allow it, generate what the reference states from the code,
such as the default values from the signature and the list of choices from
the code that validates them, rather than writing it a second time
(\DesignRef{DRY: Don't repeat yourself}).



\subsection{Document every key and every file}
\label{sec:workflow-docs}

The reference of workflow code is not its functions but what its users write
and read: the command that runs it, the configuration, and the files that go
in and come out. These form its public API (\DesignRef{Design stable APIs}),
and each part should be documented:

\begin{itemize}
  \item \textbf{Running it.} The command and its options, and the environment
    it needs (\DesignRef{Pin what you run}).
  \item \textbf{Configuration.} Every key, with its meaning, type, unit,
    default, and allowed values, and the keys that only make sense together.
  \item \textbf{Inputs.} Every input file and its format, including the
    reference data and the versions that were tested, and the table of
    samples or inputs with its columns.
  \item \textbf{Outputs.} Every output file, what it contains, its format and
    columns with their units, and which outputs are final and which are
    intermediate.
  \item \textbf{Resources.} The memory, the number of cores, the disk space,
    and the run time that each step needs, as functions of the size of the
    input (\DesignRef{Size matters}), so that users can request them on a
    cluster.
\end{itemize}

The configuration is documented best where it is defined: in a schema that
both validates the configuration (\DesignRef{Garbage in, useful error out})
and generates its reference, so that the two cannot disagree
(\DesignRef{DRY: Don't repeat yourself}). The small dataset of the
end-to-end test (\DesignRef{Test the whole workflow on a small dataset})
doubles as the dataset of a tutorial: a run that users can start right after
installing, with outputs they can compare with the documented ones.



\subsection{Comments say why, not what}
\label{sec:comments}

The code says what it does; a comment should say what the code cannot: why
it does it that way. Good comments give the reason for a choice that a
reader would otherwise undo, such as a reordering that avoids cancellation
or a loop kept for memory; the source of a formula, down to the equation
number and the sign convention of the paper; the unit or convention of a
value where its name does not say it (\DesignRef{Name axes, units, and
conventions}); and the bug report behind a workaround, so that it can be
removed once the bug is fixed. Comments that repeat the code add nothing
and fall out of date with the first change, and from then on they
contradict it.

Comments explain a few lines; some knowledge concerns the whole package.
Decisions about its design, such as why a method was chosen over the
alternatives, why a dependency was added, or why a convention differs from
the host stack's, should be recorded in the repository when they are made,
for example as short decision records, together with the conventions and
the release process. They are the documentation of the maintainers, and
without them the package depends on the memory of its authors
(\DesignRef{Code is for life, not just for Christmas}).



\subsection{Keep a changelog}
\label{sec:changelog}

\emph{Design Principles} requires a changelog that records every change of
behavior (\DesignRef{Design stable APIs}). Write it for users, not for the
developers: a list of commit messages says what was done to the code, while
users need to know what changed for them and whether they have to act. The
format of Keep a Changelog serves this well: one entry per release, with its
version and date, and the changes grouped by kind, such as added, changed,
deprecated, removed, and fixed.

Changes that can alter results should stand in a group of their own, such as
``changes of results'', rather than among the fixes, so that users can
decide at a glance whether to rerun an analysis. Each such entry says which
entry points are affected, how the results change, and how to get the old
behavior back if that is possible. Each deprecation names its replacement
and the release in which the old behavior will be removed.



\subsection{Credit where credit is due}
\label{sec:cite}

\emph{Design Principles} requires published software to provide
machine-readable citation metadata (\DesignRef{Make every result traceable
to a version}). Make citing it easy as well: users cite what they find
quickly. Software that is not cited is not credited, and a result that does
not cite it cannot be traced to it.

State how to cite the software in every place users look: the README, the
documentation, and the host stack's citation mechanism, such as the
\texttt{citation} function of R; a \texttt{CITATION.cff} file in the
repository is read by code-hosting platforms and archives. Ask for two
citations where there are two contributions, the software and the paper of
the method (Subsection~\ref{sec:paper}). Archives such as Zenodo give each
release a DOI of its own and the software a DOI for all versions: ask users
to cite the DOI of the version they ran, so that the citation also says
which version computed the result.

Software that runs the methods of others, such as a workflow that calls
established tools, should also say which of them to cite. Workflow code can
list them next to the versions of the tools that it records for each run
(\DesignRef{Make every result traceable to a version}).



% ==========================================================
\section{Keeping documentation true}
\label{sec:true}
% ==========================================================

\subsection{Docs as code}
\label{sec:docs-as-code}

Documentation should be treated like code: kept in the repository of the
package, written in plain-text formats such as Markdown, reStructuredText,
or Quarto, reviewed, built, and released with the code. Documentation kept
elsewhere, such as in a wiki, a shared drive, or a slide deck, has no
version, changes on its own schedule, and describes some version of the
software, but not a stated one. Figures that show what the software computes
should be made by code kept with the documentation, preferably by its build,
so that they show what the current version computes (\DesignRef{Keep the
recipe, not just the cake}).

A change of behavior and the change of its documentation belong in the same
commit, and the review of a change checks both. When the code and the
documentation disagree, one of them is wrong, and the disagreement is a bug
like any other: decide which describes the intended behavior, and fix the
other.

Build the documentation automatically on every change, and treat warnings,
such as broken links, references that do not resolve, and user-facing
functions without documentation, as errors. Publish it for each release, so
that readers can choose the documentation of the version they run, and show
the version on every page (\DesignRef{Make every result traceable to a
version}). The documentation of the development version may be published
too, but it should not be the only one, because it describes features that
the released version does not have.



\subsection{Every example runs}
\label{sec:examples}

Every example in the documentation should run automatically on every change:
the examples of the reference, the minimal example of the README, the
tutorials, and the how-to guides. An example that is never run breaks
silently when the code changes, and a reader who copies it gets an error at
best and a wrong result at worst. Host stacks provide the means, such as
doctests in Python and Julia, the examples and vignettes that the check of
an R package runs, and notebooks and Quarto documents executed when the
documentation is built.

Running an example shows that it runs; comparing its printed output with the
documented output, as a doctest does, shows that the output has not changed
(\TypesRef{It runs}; \TypesRef{Snapshot}). Neither shows that the result is
right, so examples do not replace tests (\DesignRef{Testing is not
optional}). Yet readers take the output that an example shows as the right
answer, so check it against an oracle before publishing it, as for any
reference output (\TestingRef{Unchanged is not correct}). Printed
floating-point numbers vary in their last digits between platforms: print
fewer digits, or use the tolerance options of the doctest tool, rather than
dropping the output.

Examples need data that are always there. Use small datasets that ship with
the package or are generated from a fixed seed. Where an example needs real
data that are too large to ship, download them from an archive with a
persistent identifier, check them against a stored checksum, and keep a
copy for the build, so that the documentation does not break when a server
moves.



\end{document}
